\def\level{BTS}  % Classe à droite
\def\course{CIEL 2}
\def\chaptername{Equations différentielles} % Titre au centre
\def\docpart{Equations différentielles~--~Exercices} % Partie du cours
\def\docname{C214}        % Identifiant à gauche

\pagestyle{cours_FP}
\makeheaderBTS{} % Affiche le bandeau

\begin{multicols}{2}

		\begin{exercice_cours}[\quad]{}

Résoudre les équations sans second membre:

\begin{enumerate}
  \item $2 y^{\prime}+3 y=0$
  \item $y^{\prime}+2 y=0$
  \item $4 y^{\prime}+5 y=0$
  \item $2 y^{\prime}-3 y=0$
\end{enumerate}

\end{exercice_cours}

		\begin{exercice_cours}[\quad]{}

Vérifier que la fonction $p$ proposée est une solution particulière, puis résoudre l'équation.

\begin{enumerate}
  \item $y^{\prime}+2 y=6$ avec $p(x)=3$.
  \item $y^{\prime}-y=x$ avec $p(x)=-x-1$.
  \item $2 y^{\prime}+y=e^{x}$ avec $p(x)=\frac{1}{3} e^{x}$.
\end{enumerate}
		\end{exercice_cours}

				\begin{exercice_cours}[\quad]{}

		Déterminer la fonction $f$ solution vérifiant la condition initiale donnée :

\begin{enumerate}
  \item $y^{\prime}-2 y=0$ et $f(0)=2$
  \item $y^{\prime}+y=0$ et $f(-1)=3$
\end{enumerate}


		\end{exercice_cours}

		\begin{exercice_cours}[\quad]{}

			On considère l'équation différentielle suivante (E) : $5 y^{\prime}-y=x$

\begin{enumerate}
  \item Vérifier que $p$ définie par $p(x)=-x-5$ est une solution de (E).
  \item En déduire toutes les solutions de (E).
  \item Déterminer la solution $f$ de (E) telle que $f(0)=1$.
\end{enumerate}


		\end{exercice_cours}

		\begin{exercice_cours}[\quad]{}

			On considère l'équation différentielle (E) : $y^{\prime}-y=x^{2}-x-1$.

\begin{enumerate}
  \item a) Écrire l'équation homogène associée.\\
b) Donner la forme générale des solutions de cette équation homogène.
  \item Vérifier que $p$ définie par $p(x)=-x^{2}-x$ est une solution particulière de (E).
  \item Déterminer toutes les solutions de (E).
  \item En déduire la solution de (E) qui vérifie $f(0)=1$.
\end{enumerate}


		\end{exercice_cours}

		\begin{exercice_cours}[\quad]{}

			On considère l'équation différentielle suivante (E) : $y^{\prime}+3 y=5$

\begin{enumerate}
  \item Déterminer le réel $a$ pour que la fonction $p$ définie sur $\mathbb{R}$ par $p(x)=a$ soit une solution de (E).
  \item Résoudre l'équation.
\end{enumerate}


		\end{exercice_cours}

		\begin{exercice_cours}[\quad]{}

			Soit l'équation (E) : $y^{\prime}+2 y=3 x+2$.

\begin{enumerate}
  \item Donner et résoudre l'équation homogène associée ( $\mathrm{E}_{0}$ ).
  \item Déterminer une solution particulière de (E) sous la forme $p(x)=a x+b$, où $a$ et $b$ sont deux réels.
  \item En déduire la solution générale de (E).
\end{enumerate}

\end{exercice_cours}



\begin{exercice_cours}[\quad]{}
On considère l'équation différentielle 

(E) : $y^{\prime}+2 y=-5 e^{-2 x}$.

\begin{enumerate}
  \item Résoudre l'équation homogène associée à (E).
  \item Montrer que la fonction $p$ définie sur $\mathbb{R}$ par $p(x)=-5 x e^{-2 x}$ est une solution particulière de (E).
  \item En déduire l'ensemble des solutions de (E).
  \item En déduire la solution $f$ de (E) qui vérifie $f(0)=1$.
\end{enumerate}

\end{exercice_cours}


		\begin{exercice_cours}[\quad]{}

			On considère l'équation différentielle 
			
			(E) : $y^{\prime}+2 y=2 x-e^{-2 x}$.

\begin{enumerate}
  \item Résoudre l'équation homogène associée à (E).
  \item Montrer que la fonction $p$ définie sur $\mathbb{R}$ par $p(x)=x-\dfrac{1}{2}-x e^{-2 x}$ est une solution particulière de (E).
  \item En déduire l'ensemble des solutions de (E).
  \item En déduire la solution $f$ de (E) qui vérifie $f(0)=0$.
\end{enumerate}


\end{exercice_cours}


\end{multicols}